% Generated draft: requires mathematical and source review.
\documentclass[11pt]{article}
\usepackage[margin=25mm]{geometry}
\usepackage{fontspec}
\setmainfont{Latin Modern Roman}
\usepackage{amsmath,amssymb,mathtools,unicode-math}
\setmathfont{Latin Modern Math}
\usepackage{enumitem,needspace,xurl,hyperref}
\setlength{\parindent}{0pt}
\setlength{\parskip}{5pt}
\setlength{\emergencystretch}{3em}
\newcommand{\N}{\mathbb{N}}
\newcommand{\Z}{\mathbb{Z}}
\newcommand{\Q}{\mathbb{Q}}
\newcommand{\R}{\mathbb{R}}
\newcommand{\C}{\mathbb{C}}
\newcommand{\F}{\mathbb{F}}
\newcommand{\A}{\mathbb{A}}
\newcommand{\PP}{\mathbb{P}}
\newcommand{\E}{\mathbb{E}}
\newcommand{\eps}{\varepsilon}
\newcommand{\dd}{\,\mathrm{d}}
\begin{document}
\section*{UnsolvedMath ID: 800753}
\textbf{Source problem code:} catalogue entry
\subsection{Definitions and mathematical statement}
For the standard real-scalar formulation, let \(X\) be a real Banach space and let \(X^*\) denote its space of continuous real-valued linear functionals. A \emph{support set} in \(X\) is a nonempty convex subset \(C\subseteq X\) such that
\[
\forall x\in C\;\exists f\in X^*:\qquad
f(x)=\inf_{y\in C}f(y)<\sup_{y\in C}f(y).
\]
Thus the functional may depend on \(x\); it attains its minimum over \(C\) at \(x\) and is not constant on \(C\). The supremum may be \(+\infty\). No closedness or boundedness assumption on \(C\) is imposed in the supplied definition.

\textbf{Problem (S. Rolewicz, as attributed in the supplied source).} Characterize the Banach spaces that admit support sets. In particular, is it true that every nonseparable real Banach space \(X\) contains such a set \(C\)? Here nonseparability means that \(X\) has no countable norm-dense subset.

\textbf{Formulation qualification.} The source does not explicitly specify the scalar field or continuity of the functionals. The formulation above uses continuous linear functionals and real scalars, as conventionally understood in this ordered support-functional definition. For complex Banach spaces, an additional convention, such as passing to the underlying real space or using real parts of complex functionals, would be needed; the source does not state one.

\textbf{Supplied background.} The source attributes the problem to S. Rolewicz and reports that he showed that separable Banach spaces admit no support sets and observed that many nonseparable Banach spaces do admit them. These are supplied historical assertions, not independently verified conclusions of this draft.
\subsection{Short English statement}
A support set is a nonempty convex set for which every point minimizes some continuous linear functional that is nonconstant on the set. The problem, attributed to S. Rolewicz, asks which Banach spaces contain support sets and whether every nonseparable Banach space does. The displayed formulation uses real scalars; the source leaves the scalar-field convention implicit.
\subsection{Sources}
\begin{itemize}
\item Local UnsolvedMath registry, permanent ID \(800753\), exact source problem code \texttt{catalogue entry}; title: \emph{Banach Spaces and Ramsey Theory --- Problem 17}. The registry labels the problem \texttt{open}; this is a supplied status label, not an independently verified current assessment.
\item Dodos and Lopez-Abad, \emph{Banach Spaces and Ramsey Theory: Some Open Problems} (2010), Problem 17. Supplied source URL: \url{https://arxiv.org/abs/1006.2668}. The record reports access on 2026-10-16 and extraction by \texttt{arxiv\_latex}; these are supplied provenance details, not a claim of access or verification by this draft.
\item S. Rolewicz, reference label [Rol], as cited and credited in the supplied problem text. The record supplies no full bibliographic entry for [Rol], so the original reference and the attributed background results cannot be identified precisely from this record alone.
\end{itemize}
\end{document}
